\documentclass[12pt,a4paper]{article}
\usepackage[utf8]{inputenc}
\usepackage[T1]{fontenc}
\usepackage{geometry}
\usepackage{amsmath}
\usepackage{amssymb}
\usepackage{listings}
\usepackage{xcolor}
\usepackage{hyperref}
\usepackage{booktabs}
\usepackage{longtable}
\usepackage{array}
\usepackage{enumitem}
\usepackage{caption}
\usepackage{fancyhdr}
\usepackage{titlesec}
\usepackage{setspace}
\usepackage{microtype}
\usepackage{verbatim}
\usepackage{multirow}
\usepackage{float}
\usepackage[bottom]{footmisc}

\geometry{a4paper, margin=1in, headheight=14.5pt, footskip=30pt, textheight=24cm}

\definecolor{codegreen}{rgb}{0,0.6,0}
\definecolor{codegray}{rgb}{0.5,0.5,0.5}
\definecolor{codeblue}{rgb}{0,0,0.95}
\definecolor{backcolour}{rgb}{0.95,0.95,0.92}
\definecolor{headerblue}{rgb}{0.0,0.3,0.6}
\definecolor{footerblue}{rgb}{0.2,0.4,0.7}
\definecolor{warningbg}{rgb}{1.0,0.95,0.9}
\definecolor{passgreen}{rgb}{0.0,0.6,0.0}

\lstdefinestyle{mystyle}{
    backgroundcolor=\color{backcolour},
    commentstyle=\color{codegreen},
    keywordstyle=\color{codeblue}\bfseries,
    numberstyle=\tiny\color{codegray},
    stringstyle=\color{codegreen},
    basicstyle=\ttfamily\footnotesize,
    breakatwhitespace=false,
    breaklines=true,
    captionpos=b,
    keepspaces=true,
    numbers=left,
    numbersep=5pt,
    showspaces=false,
    showstringspaces=false,
    showtabs=false,
    tabsize=2,
    frame=single,
    frameround=tttt
}

\lstdefinestyle{outputstyle}{
    backgroundcolor=\color{backcolour},
    basicstyle=\ttfamily\footnotesize\color{codegray},
    breakatwhitespace=false,
    breaklines=true,
    captionpos=b,
    keepspaces=true,
    numbers=none,
    showspaces=false,
    showstringspaces=false,
    showtabs=false,
    tabsize=2,
    frame=single,
    frameround=tttt
}

\lstset{style=mystyle}

\pagestyle{fancy}
\fancyhf{}
\fancyhead[L]{\textcolor{headerblue}{\textbf{SiliconeSignal Technologies}}}
\fancyhead[R]{\textcolor{headerblue}{\textbf{Gemmini Backend Reference Manual}}}
\fancyfoot[L]{\textcolor{footerblue}{\small Meknes, Morocco}}
\fancyfoot[C]{\thepage}
\fancyfoot[R]{\textcolor{footerblue}{\small Version 3.1}}

\titleformat{\section}
    {\normalfont\Large\bfseries\color{headerblue}}
    {\thesection}{1em}{}

\titleformat{\subsection}
    {\normalfont\large\bfseries\color{headerblue}}
    {\thesubsection}{1em}{}

\titleformat{\subsubsection}
    {\normalfont\bfseries\color{headerblue}}
    {\thesubsubsection}{1em}{}

\newenvironment{definition}[2][Definition]{
    \begin{center}
    \begin{tabular}{|p{0.95\textwidth}|}
    \hline
    \textbf{#1: #2} \\
    \hline
    \begin{minipage}{0.93\textwidth}
}{
    \end{minipage} \\
    \hline
    \end{tabular}
    \end{center}
}

\newenvironment{notebox}[1][Note]{
    \begin{center}
    \begin{tabular}{|p{0.95\textwidth}|}
    \hline
    \textbf{#1} \\
    \hline
    \begin{minipage}{0.93\textwidth}
}{
    \end{minipage} \\
    \hline
    \end{tabular}
    \end{center}
}

\newenvironment{warningbox}[1][Warning]{
    \begin{center}
    \begin{tabular}{|p{0.95\textwidth}|}
    \hline
    \textcolor{red}{\textbf{#1}} \\
    \hline
    \begin{minipage}{0.93\textwidth}
}{
    \end{minipage} \\
    \hline
    \end{tabular}
    \end{center}
}

\newcommand{\code}[1]{\texttt{#1}}
\newcommand{\gemmini}{\textsc{Gemmini}}
\newcommand{\onnxmlir}{\texttt{ONNX-MLIR}}

\title{
    \vspace{2cm}
    {\Huge\bfseries\color{headerblue}Gemmini Accelerator Backend} \\
    {\Large Reference Manual} \\
    \vspace{0.5cm}
    {\large ONNX-MLIR Integration \& Implementation Guide} \\
    \vspace{1cm}
    {\large Version 3.1 (Corrected)}
}
\author{\textbf{SST Team}}
\date{}

\begin{document}

\maketitle
\thispagestyle{empty}
\newpage

\section*{Document Control}

\begin{tabular}{|p{0.25\textwidth}|p{0.7\textwidth}|}
\hline
\textbf{Document ID} & SST-GEMMINI-RM-008 \\
\hline
\textbf{Version} & 3.1 \\
\hline
\textbf{Status} & Final \\
\hline
\textbf{Classification} & Technical Reference \\
\hline
\textbf{Author} & SST Team \\
\hline
\textbf{Last Updated} & \today \\
\hline
\textbf{Approval} & Pending \\
\hline
\end{tabular}

\vspace{1cm}

\noindent\textbf{Abstract:} This reference manual provides comprehensive documentation for the Gemmini accelerator backend within the ONNX-MLIR compiler framework. The backend enables ONNX models to target the Gemmini systolic array accelerator on RISC-V platforms. As of June 2026, the backend passes 41/41 MLIR lit tests, 18/18 integration tests, 22/22 example tests, and 9/9 model zoo models. This document covers architecture, compiler passes, MLIR dialects, runtime library, build configuration, testing infrastructure, validation results, usage examples, and troubleshooting.

\newpage

\tableofcontents
\newpage

\section{Introduction}

\subsection{Document Purpose and Scope}

This reference manual serves as the authoritative guide for developers, researchers, and system architects working with the Gemmini accelerator backend in ONNX-MLIR. It documents the complete implementation, from high-level architecture to low-level runtime details.

\subsection{Core Definitions}

\begin{definition}{Gemmini}
Gemmini is an open-source, configurable systolic array generator developed by UC Berkeley. It produces RISC-V RoCC (RISC-V Custom Coprocessor) accelerators for dense linear algebra operations, particularly matrix multiplication and convolution. The generator is written in Chisel and produces synthesizable Verilog for FPGA or ASIC implementations.

\textbf{Key properties:}
\begin{itemize}[leftmargin=*]
    \item Configurable systolic array dimensions (default 16x16)
    \item Support for weight-stationary (WS) and output-stationary (OS) dataflows
    \item Software-managed scratchpad memory
    \item DMA engine for data movement
    \item Optional activation functions (ReLU, ReLU6)
\end{itemize}
\end{definition}

\begin{definition}{Systolic Array}
A systolic array is a homogeneous network of processing elements (PEs) that rhythmically compute and pass data through the network. Gemmini's default configuration uses a 16x16 systolic array where each PE performs a multiply-accumulate (MAC) operation every cycle.

\textbf{Characteristics:}
\begin{itemize}[leftmargin=*]
    \item Data flows through the array in a rhythmic, pipelined fashion
    \item Each PE communicates only with its nearest neighbors
    \item Reduces memory bandwidth requirements by reusing data
    \item Achieves high throughput for matrix multiplication
\end{itemize}
\end{definition}

\begin{definition}{RoCC (RISC-V Custom Coprocessor)}
RoCC is a RISC-V extension that defines a standard interface for attaching custom coprocessors. Gemmini is implemented as a RoCC accelerator using the RoCC port of a Rocket or BOOM tile.

\textbf{Instruction format:}
\begin{lstlisting}[language=C]
.insn r CUSTOM_3, 0x3, funct7, x0, rs1, rs2
\end{lstlisting}
\end{definition}

\begin{definition}{Scratchpad Memory}
Software-managed on-chip SRAM, distinct from hardware-managed caches. Gemmini's scratchpad has 16,384 rows x 16 bytes per row (256 KB total) organized in 4 banks for parallel access. This is 16x larger than the stub values used in early development.
\end{definition}

\begin{definition}{Double Buffering (Ping-Pong)}
A technique that uses two memory buffers to overlap data transfer with computation. While the systolic array computes on one buffer, DMA loads the next tile into the other buffer. The slot offsets are expressed as relative row ranges within the 16,384-row address space.
\end{definition}

\begin{definition}{Spike Simulator}
The official RISC-V ISA simulator. With the --extension=gemmini flag, Spike loads a shared library (libgemmini.so) that emulates Gemmini RoCC instructions.
\end{definition}

\begin{definition}{ONNX-MLIR}
An open-source compiler that translates ONNX (Open Neural Network Exchange) models into MLIR (Multi-Level Intermediate Representation) and then to LLVM IR, ultimately producing executable binaries for various targets including RISC-V.
\end{definition}

\subsection{Current Status Summary}

\begin{table}[H]
\centering
\begin{tabular}{|p{0.4\textwidth}|p{0.2\textwidth}|p{0.3\textwidth}|}
\hline
\textbf{Test Suite} & \textbf{Count} & \textbf{Status} \\
\hline
MLIR lit tests & 41 & \textcolor{passgreen}{41/41 PASS} \\
\hline
Integration tests & 18 & \textcolor{passgreen}{18/18 PASS} \\
\hline
Example tests & 22 & \textcolor{passgreen}{22/22 PASS} \\
\hline
Model zoo & 9 & \textcolor{passgreen}{9/9 PASS} \\
\hline
Numerical correctness (Spike) & 3 models & \textcolor{passgreen}{PASS at 1e-5} \\
\hline
\end{tabular}
\caption{Current Validation Status}
\end{table}

\section{Quick Start Guide (5 Minutes)}

\subsection{Step 1: Check Prerequisites}

\begin{lstlisting}[language=bash]
scripts/check_prerequisites.sh --mode compile-only
\end{lstlisting}

\subsection{Step 2: Build the Compiler}

\begin{lstlisting}[language=bash]
cmake -S . -B gemmini_toolchain_build -DONNX_MLIR_ENABLE_GEMMINI=ON -DONNX_MLIR_GEMMINI_MODE=compile-only -DCMAKE_BUILD_TYPE=Release
cmake --build gemmini_toolchain_build -j$(nproc) --target onnx-mlir
\end{lstlisting}

\subsection{Step 3: Compile ResNet-18}

\begin{lstlisting}[language=bash]
ONNX_MLIR=gemmini_toolchain_build/Release/bin/onnx-mlir
$ONNX_MLIR --maccel=Gemmini --EmitMLIR resnet18-v1-7.onnx -o /tmp/resnet18_gemmini
\end{lstlisting}

\subsection{Step 4: Count Gemmini Calls}

\begin{lstlisting}[language=bash]
grep -c 'om_gemmini_' /tmp/resnet18_gemmini.onnx.mlir
\end{lstlisting}

\begin{lstlisting}[style=outputstyle]
21
\end{lstlisting}

\subsection{Step 5: Run on Spike (Simulator Mode)}

\begin{lstlisting}[language=bash]
RISCV_INSTALL="$HOME/riscv-gemmini" scripts/setup_gemmini_simulator.sh --install
cmake -S . -B gemmini_toolchain_build_sim -DONNX_MLIR_ENABLE_GEMMINI=ON -DONNX_MLIR_GEMMINI_MODE=simulator -DRISCV_INSTALL="$HOME/riscv-gemmini"
cmake --build gemmini_toolchain_build_sim -j$(nproc)
$ONNX_MLIR --maccel=Gemmini --mtriple=riscv64-unknown-linux-gnu resnet18-v1-7.onnx -o resnet18_rv64
spike --extension=gemmini --isa=rv64imafdc "$HOME/riscv-gemmini/bin/pk" ./resnet18_rv64
\end{lstlisting}

\subsection{Expected Output}

\begin{table}[H]
\centering
\begin{tabular}{|p{0.4\textwidth}|p{0.55\textwidth}|}
\hline
\textbf{Item} & \textbf{Expected Value} \\
\hline
Gemmini runtime calls & 21 (20 om\_gemmini\_conv\_f32 + 1 om\_gemmini\_gemm\_f32\_bias) \\
\hline
MLIR lit tests pass rate & 41/41 (100\%) \\
\hline
Integration tests pass rate & 18/18 (100\%) \\
\hline
\end{tabular}
\caption{Expected Quick Start Results}
\end{table}

\section{Architecture Overview}

\subsection{Compilation Pipeline}

\begin{verbatim}
ONNX Model (.onnx)
       |
       v
ONNX Dialect
       |
       v
ONNXToGemmini Pass (benefit=10)
       |
   +---+---+
   |       |
   v       v
Gemmini   krnl.call
Dialect   (Runtime Delegate)
   |
   v
GemminiTiling
   |
   v
GemminiToGemminiLow
   |
   v
GemminiLowRewrite
   |
   v
GemminiLowToLLVM
   |
   v
LLVM Dialect -> RISC-V Binary
\end{verbatim}

\subsection{Two Execution Paths}

\begin{table}[H]
\centering
\begin{tabular}{|p{0.25\textwidth}|p{0.35\textwidth}|p{0.35\textwidth}|}
\hline
\textbf{Aspect} & \textbf{Runtime Delegate Path} & \textbf{Direct RoCC Path} \\
\hline
Primary ops & Conv, Gemm, MatMul f32, Resize, Pad, Concat, Slice, Transpose, Split, Sigmoid, Mul, Relu, BN, Pooling, Softmax & MatMul int8/i32 \\
\hline
Lowering & krnl.call @om\_gemmini\_* & gemmini\_low.* to CUSTOM\_3 asm \\
\hline
Hardware use & Quantize f32 to int8, call tiled\_matmul\_auto, dequantize & Direct RoCC instructions \\
\hline
Benefit & 10 (wins over standard path) & 11 (highest priority) \\
\hline
\end{tabular}
\caption{Two Execution Paths Comparison}
\end{table}

\subsection{Hardware Parameters}

\begin{table}[H]
\centering
\begin{tabular}{|p{0.35\textwidth}|p{0.25\textwidth}|p{0.35\textwidth}|}
\hline
\textbf{Parameter} & \textbf{Value} & \textbf{Description} \\
\hline
Systolic array dimension & DIM = 16 & 16x16 processing elements \\
\hline
Scratchpad banks & 4 & Independent memory banks \\
\hline
Scratchpad rows & 16,384 & Total addressable rows (256 KB) \\
\hline
Accumulator rows & 1,024 & Accumulator buffer rows (64 KB) \\
\hline
Element type & int8 & Input element type \\
\hline
Accumulator type & int32 & Accumulator element type \\
\hline
RoCC channel & XCUSTOM\_ACC = 3 & CUSTOM\_3 opcode (0x7b) \\
\hline
Dataflow mode & Weight-Stationary (WS) & Gemmini architecture \\
\hline
\end{tabular}
\caption{Gemmini Hardware Parameters}
\end{table}

\section{Compiler Pipeline Reference}

\subsection{Lowering Passes}

\begin{table}[H]
\centering
\begin{tabular}{|p{0.3\textwidth}|p{0.3\textwidth}|p{0.35\textwidth}|}
\hline
\textbf{Pass} & \textbf{File Location} & \textbf{Role} \\
\hline
ONNXToGemmini & Conversion/ONNXToGemmini/ & Intercept supported ONNX ops; emit gemmini dialect or krnl.call \\
\hline
GemminiTiling & Transform/Gemmini/ & Tile direct MatMul to DIM=16 blocks; assign ping-pong scratchpad slots \\
\hline
GemminiToGemminiLow & Transform/GemminiLow/ & Convert high-level gemmini ops to instruction-level gemmini\_low ops \\
\hline
GemminiLowRewrite & Transform/GemminiLow/ & Peephole: eliminate redundant adjacent config/fence/memory ops \\
\hline
StaticScratchpadAllocation & Transform/GemminiLow/ & Assign fixed scratchpad row offsets to each tile \\
\hline
GemminiLowToLLVM & Conversion/GemminiLowToLLVM/ & Emit RISC-V CUSTOM\_3 inline assembly \\
\hline
\end{tabular}
\caption{Gemmini Lowering Passes}
\end{table}

\subsection{Corrected RoCC Instruction Sequence}

The 10-step sequence below is the exact order emitted by GemminiLowToLLVM.cpp for a 16x16 WS-mode direct MatMul. Note that FENCE is emitted before every MVOUT.

\begin{table}[H]
\centering
\begin{tabular}{|p{0.25\textwidth}|p{0.35\textwidth}|p{0.35\textwidth}|}
\hline
\textbf{Step} & \textbf{Instruction} & \textbf{Key Operands} \\
\hline
1 & CONFIG\_EX (funct=0) & rs1 = (acc\_scale << 32) | (1 << 16) | dataflow\_WS \\
\hline
2 & CONFIG\_MVIN\_A (funct=0) & rs1 = (MVIN\_SCALE\_ID << 32) | 0x01, rs2 = stride\_bytes \\
\hline
3 & MVIN\_A (funct=2) & rs2 = (rows<<48)|(cols<<32)|spad\_row\_A \\
\hline
4 & CONFIG\_MVIN\_B (funct=0) & Same as Step 2 \\
\hline
5 & MVIN\_B (funct=2) & rs2 = (rows<<48)|(cols<<32)|spad\_row\_B \\
\hline
6 & PRELOAD (funct=6) & rs2 = C\_acc\_addr | 0x80000000 \\
\hline
7 & COMPUTE\_PRELOADED (funct=4) & rs2 = 0xFFFFFFFF (no-bias sentinel) \\
\hline
8 & FENCE & Memory and pipeline flush \\
\hline
9 & CONFIG\_MVOUT (funct=0) & rs1 = 0x02, rs2 = (acc\_scale<<32) | stride\_bytes \\
\hline
10 & MVOUT (funct=3) & rs2 = (rows<<48)|(cols<<32)|acc\_offset|0x80000000 \\
\hline
\end{tabular}
\caption{Corrected RoCC Instruction Sequence (10 steps with FENCE)}
\end{table}

\subsection{Cycle Counting Formula}

The total cycle count for a 16x16 WS-mode direct MatMul is:

\[
\begin{aligned}
\text{Total Cycles} = & \underbrace{1}_{\text{CONFIG\_EX}} + \underbrace{1}_{\text{CONFIG\_MVIN\_A}} + \underbrace{64}_{\text{MVIN\_A}} + \underbrace{1}_{\text{CONFIG\_MVIN\_B}} + \underbrace{64}_{\text{MVIN\_B}} \\
& + \underbrace{1}_{\text{PRELOAD}} + \underbrace{256}_{\text{COMPUTE\_PRELOADED}} + \underbrace{1}_{\text{FENCE}} + \underbrace{1}_{\text{CONFIG\_MVOUT}} + \underbrace{64}_{\text{MVOUT}} \\
= & 454 \text{ cycles}
\end{aligned}
\]

% FIX: title was "Six Correctness Invariants" but lists 7 items — corrected to Seven.
\begin{warningbox}[Seven Correctness Invariants]
\begin{enumerate}[leftmargin=*]
    \item Use CUSTOM\_3 (opcode 0x7b), not CUSTOM\_0
    \item CONFIG\_MVIN required before each MVIN
    \item PRELOAD required before COMPUTE\_PRELOADED
    \item rs2 = 0xFFFFFFFF in COMPUTE\_PRELOADED (no-bias sentinel)
    \item Bit 31 set in MVOUT rs2 (read from accumulator)
    \item FENCE required before CONFIG\_MVOUT and MVOUT
    \item a\_stride=1 (bit 16 of CONFIG\_EX rs1), c\_stride=1 (bit 48 of rs2)
\end{enumerate}
\end{warningbox}

\section{Supported Operators}

\subsection{Complete Operator Table}

\begin{longtable}{|p{0.28\textwidth}|p{0.2\textwidth}|p{0.47\textwidth}|}
\hline
\textbf{ONNX Op} & \textbf{Path} & \textbf{Runtime Function / Constraints} \\
\hline
\endfirsthead
\hline
\textbf{ONNX Op} & \textbf{Path} & \textbf{Runtime Function / Constraints} \\
\hline
\endhead
MatMul int8 & Direct RoCC & gemmini\_low.* (no call); rank-2 static, 16x16 WS tiles \\
\hline
MatMul f32 & Runtime & om\_gemmini\_matmul\_f32\_hw; rank-2 static; quantize to Gemmini to dequantize \\
\hline
MatMul f32 batched & Runtime & om\_gemmini\_matmul\_f32\_nd\_hw; rank-N x rank-2 \\
\hline
MatMulInteger & Runtime & om\_gemmini\_matmulinteger\_i8i8acc32[\_zp]; i8xi8 to i32 (integration test coverage) \\
\hline
QLinearConv & Runtime & om\_gemmini\_qlinearconv\_i8[\_bias]; rank-4 int8 NCHW; batch=1; group=1 (integration test coverage) \\
\hline
Conv f32 & Runtime & om\_gemmini\_conv\_f32[\_bias]; rank-4 NCHW; dynamic batch; static C/H/W \\
\hline
ConvTranspose f32 & Runtime & om\_gemmini\_convtranspose\_f32[\_bias]; col2im scatter-add; dynamic batch \\
\hline
Gemm f32 & Runtime & om\_gemmini\_gemm\_f32[\_bias]; rank-2; transA/transB/alpha/beta \\
\hline
Relu f32 & Runtime & om\_gemmini\_relu\_f32; rank-4; dynamic batch \\
\hline
BatchNormalization f32 & Runtime & om\_gemmini\_batchnorm\_f32; rank-4 NCHW; inference only \\
\hline
Add f32 & Runtime & om\_gemmini\_add\_f32; rank-4 same-shape; dynamic batch \\
\hline
MaxPool f32 & Runtime & om\_gemmini\_maxpool\_f32; rank-4 NCHW; static kernel/stride/pads \\
\hline
AveragePool f32 & Runtime & om\_gemmini\_avgpool\_f32; rank-4 NCHW; static kernel/stride/pads \\
\hline
GlobalAveragePool f32 & Runtime & om\_gemmini\_globalavgpool\_f32; rank-4 NCHW \\
\hline
Softmax f32 & Runtime & om\_gemmini\_softmax\_f32; rank-2 or rank-4 \\
\hline
Resize f32 & Runtime & om\_gemmini\_resize\_nearest\_f32 / om\_gemmini\_resize\_linear\_f32 \\
\hline
Pad f32 & Runtime & om\_gemmini\_pad\_constant/reflect/edge\_f32; static spatial pads \\
\hline
Concat f32 & Runtime & om\_gemmini\_concat\_f32; 2-input rank-4 NCHW; any axis \\
\hline
Slice f32 & Runtime & om\_gemmini\_slice\_f32; rank-4; constant axes/starts/ends/steps \\
\hline
Transpose f32 & Runtime & om\_gemmini\_transpose\_f32\_hw; rank-4; constant perm \\
\hline
Split f32 & Runtime & om\_gemmini\_split\_2/3/4\_f32; rank-4; 2-4 outputs; any axis \\
\hline
Sigmoid f32 & Runtime & om\_gemmini\_sigmoid\_f32; rank-4; elementwise; dynamic batch \\
\hline
Mul f32 & Runtime & om\_gemmini\_mul\_f32; rank-4 same-shape; dynamic batch \\
\hline
\end{longtable}

\section{Runtime Library Reference}

\subsection{Runtime Function Groups}

\begin{table}[H]
\centering
\begin{tabular}{|p{0.25\textwidth}|p{0.3\textwidth}|p{0.4\textwidth}|}
\hline
\textbf{Group} & \textbf{Count} & \textbf{Example Symbols} \\
\hline
Integer / quantized & 4 & om\_gemmini\_matmulinteger\_i8i8acc32[\_zp], om\_gemmini\_qlinearconv\_i8[\_bias] \\
\hline
Float HW path & 7 & om\_gemmini\_matmul\_f32\_hw, om\_gemmini\_conv\_f32[\_bias], om\_gemmini\_gemm\_f32[\_bias] \\
\hline
Elementwise / norm & 6 & om\_gemmini\_relu\_f32, om\_gemmini\_batchnorm\_f32, om\_gemmini\_sigmoid\_f32, om\_gemmini\_mul\_f32 \\
\hline
Pooling & 3 & om\_gemmini\_globalavgpool\_f32, om\_gemmini\_maxpool\_f32, om\_gemmini\_avgpool\_f32 \\
\hline
Spatial & 11 & om\_gemmini\_resize\_nearest/linear\_f32, om\_gemmini\_pad\_constant/reflect/edge\_f32, om\_gemmini\_concat\_f32, om\_gemmini\_slice\_f32, om\_gemmini\_split\_2/3/4\_f32, om\_gemmini\_transpose\_f32\_hw \\
\hline
Linear algebra & 2 & om\_gemmini\_gemm\_f32, om\_gemmini\_gemm\_f32\_bias \\
\hline
\end{tabular}
\caption{Runtime Function Groups (33 total public symbols)}
\end{table}

\subsection{OMTensor Data Structure}

\begin{lstlisting}[caption=OMTensor Structure, language=C]
typedef struct OMTensor {
    void *data;           // Pointer to tensor data (row-major)
    int64_t *shape;       // Array of dimension sizes
    int64_t *strides;     // Array of strides in elements
    int64_t rank;         // Number of dimensions
    int64_t numElements;  // Total elements (product of shape)
    onnxDataType dataType;// ONNX data type enum
} OMTensor;
\end{lstlisting}

\section{Build System Reference}

\subsection{CMake Configuration Variables}

\begin{table}[H]
\centering
\begin{tabular}{|p{0.35\textwidth}|p{0.6\textwidth}|}
\hline
\textbf{Variable} & \textbf{Description} \\
\hline
ONNX\_MLIR\_ENABLE\_GEMMINI & Master switch (ON/OFF). Appends Gemmini to ONNX\_MLIR\_ACCELERATORS. \\
\hline
ONNX\_MLIR\_GEMMINI\_MODE & Build mode: compile-only, hardware, or simulator \\
\hline
GEMMINI\_ROOT & Path to Gemmini installation containing gemmini.h (default: ./gemmini) \\
\hline
RISCV\_INSTALL & RISC-V toolchain and simulator prefix \\
\hline
RISCV\_GCC & RISC-V cross-compiler command (default: riscv64-linux-gnu-gcc) \\
\hline
\end{tabular}
\caption{CMake Configuration Variables}
\end{table}

\subsection{Build Commands}

\begin{lstlisting}[caption=Compile-Only Build, language=bash]
cmake -S . -B gemmini_toolchain_build \
    -DONNX_MLIR_ENABLE_GEMMINI=ON \
    -DONNX_MLIR_GEMMINI_MODE=compile-only \
    -DONNX_MLIR_BUILD_TESTS=ON \
    -DCMAKE_BUILD_TYPE=Release \
    -G Ninja

cmake --build gemmini_toolchain_build -j"$(nproc)"
\end{lstlisting}

\begin{lstlisting}[caption=Simulator Build, language=bash]
cmake -S . -B gemmini_toolchain_build \
    -DONNX_MLIR_ENABLE_GEMMINI=ON \
    -DONNX_MLIR_GEMMINI_MODE=simulator \
    -DRISCV_INSTALL="$HOME/riscv-gemmini" \
    -DONNX_MLIR_BUILD_TESTS=ON \
    -G Ninja
\end{lstlisting}

\subsection{Prerequisite Checking}

\begin{lstlisting}[caption=Prerequisite Check Examples, language=bash]
# Check compile-only dependencies
scripts/check_prerequisites.sh --mode compile-only

# Check hardware toolchain
scripts/check_prerequisites.sh --mode hardware

# Check simulator installation
RISCV_INSTALL=/path/to/riscv ./scripts/check_prerequisites.sh --mode simulator
\end{lstlisting}

\section{Testing Infrastructure}

\subsection{Test Hierarchy}

% FIX: "check-gemmini-lit" is not a real CMake target; corrected to the actual
% add_lit_testsuite name check-onnx_mlir-accelerators-gemmini.
\begin{table}[H]
\centering
\begin{tabular}{|p{0.35\textwidth}|p{0.3\textwidth}|p{0.3\textwidth}|}
\hline
\textbf{Suite} & \textbf{Count} & \textbf{Command} \\
\hline
MLIR lit tests & 41 & check-onnx\_mlir-accelerators-gemmini \\
\hline
Integration tests & 18 & check-gemmini-integration \\
\hline
Example tests & 22 & bash gemmini/tests/run\_all\_tests.sh \\
\hline
Model zoo & 9 & run\_model\_zoo.py \\
\hline
Numerical correctness & 3 models & simulator\_numerical\_correctness.py \\
\hline
\end{tabular}
\caption{Test Suite Overview}
\end{table}

\subsection{MLIR Lit Test Categories}

\begin{table}[H]
\centering
\begin{tabular}{|p{0.3\textwidth}|p{0.65\textwidth}|}
\hline
\textbf{Category} & \textbf{Tests} \\
\hline
MatMul & matmul\_fp32.mlir, matmul\_fp32\_batched.mlir, matmul\_i8.mlir, matmul\_fp16.mlir \\
\hline
Direct RoCC & direct\_matmul\_rocc.mlir, direct\_matmul\_rocc\_edge.mlir, direct\_matmul\_i8\_rocc.mlir, config\_transpose\_bits.mlir \\
\hline
Conv / Gemm & conv2d\_gemmini.mlir, conv2d\_dynamic\_batch.mlir, conv2d\_dynamic\_spatial.mlir, convtranspose\_gemmini.mlir, gemm\_gemmini.mlir, gemm\_dynamic\_batch.mlir \\
\hline
INT8 (integration coverage) & QLinearConv and MatMulInteger zero-point variants covered by integration tests (not lit) \\
\hline
Resize & resize\_nearest.mlir, resize\_linear.mlir \\
\hline
Pad & pad\_constant.mlir, pad\_reflect.mlir, pad\_edge.mlir \\
\hline
Concat & concat\_axis0.mlir, concat\_axis1.mlir, concat\_axis2.mlir, concat\_dynamic\_batch.mlir \\
\hline
Slice & slice\_start\_end.mlir, slice\_negative\_indices.mlir, slice\_step.mlir, slice\_dynamic\_batch.mlir \\
\hline
Transpose & transpose\_identity.mlir, transpose\_nchw\_to\_nhwc.mlir, transpose\_swap\_hw.mlir, transpose\_dynamic\_batch.mlir \\
\hline
Split & split\_axis1.mlir, split\_axis2.mlir, split\_dynamic\_batch.mlir, split\_3\_outputs.mlir \\
\hline
Elementwise & sigmoid\_static.mlir, sigmoid\_dynamic\_batch.mlir, mul\_static.mlir, mul\_dynamic\_batch.mlir \\
\hline
Infrastructure & gemmini\_low\_rewrite.mlir, gemmini\_double\_buffer.mlir \\
\hline
\end{tabular}
\caption{MLIR Lit Test Categories (41 total)}
\end{table}

\subsection{Example Test Tiers}

\begin{table}[H]
\centering
\begin{tabular}{|p{0.15\textwidth}|p{0.2\textwidth}|p{0.6\textwidth}|}
\hline
\textbf{Tier} & \textbf{Tests} & \textbf{Description} \\
\hline
1 & 01, 11-14, 17, 19-transpose, 20-split, 21-sigmoid, 22-mul & f32 smoke; generate model to compile to check symbol \\
\hline
2 & 03, 04, 05, 07, 08 & INT8: MatMulInteger (scalar/vector zp), QLinearConv, direct-RoCC MatMul \\
\hline
3 & 09, 10, 15, 16, 18 & Multi-op networks; compile + link to RV64 ELF \\
\hline
4 & 06 (ResNet-18), 19-yolo-real (YOLOX-nano) & External models; SKIP gracefully when absent \\
\hline
\end{tabular}
\caption{Example Test Tiers (22 total)}
\end{table}

\subsection{Model Zoo Results}

\begin{table}[H]
\centering
\begin{tabular}{|p{0.25\textwidth}|p{0.15\textwidth}|p{0.15\textwidth}|p{0.15\textwidth}|p{0.15\textwidth}|}
\hline
\textbf{Model} & \textbf{Category} & \textbf{Size} & \textbf{Gemmini ops} & \textbf{Result} \\
\hline
ResNet-18 v1 & Classification & 44.7 MB & 21 & PASS \\
\hline
ResNet-50 v1 & Classification & 98 MB & 54 & PASS \\
\hline
MobileNetV2 & Classification & 13.3 MB & 36 & PASS \\
\hline
SqueezeNet 1.1 & Classification & 4.7 MB & 64 & PASS \\
\hline
DenseNet-121 & Classification & 32 MB & 426 & PASS \\
\hline
Res2Net-101 & Classification & 173 MB & 467 & PASS \\
\hline
YOLOX-nano & Detection & 3.5 MB & 332 & PASS + Spike \\
\hline
YOLOv5s (FP32) & Detection & 28 MB & 201 & PASS \\
\hline
BERT-tiny & NLP & 43 MB & 24 & PASS \\
\hline
\end{tabular}
\caption{Model Zoo Results (9/9 PASS)}
\end{table}

\section{Performance Characterization}

\subsection{RoCC Instruction Latency}

% FIX: original had "\hlineCONFIG_MVIN" (merged into one token); split into
% separate \hline and row entries so LaTeX can parse the table correctly.
\begin{table}[H]
\centering
\begin{tabular}{|p{0.35\textwidth}|p{0.2\textwidth}|p{0.4\textwidth}|}
\hline
\textbf{Instruction} & \textbf{Latency (cycles)} & \textbf{Notes} \\
\hline
CONFIG\_EX & 1 & Configuration register write \\
\hline
CONFIG\_MVIN & 1 & Load stride configuration \\
\hline
MVIN (per 16x16 tile) & 64 & 16x16 tile, 4 banks \\
\hline
PRELOAD & 1 & Register setup \\
\hline
COMPUTE\_PRELOADED (16x16) & 256 & DIM\textsuperscript{2} = 256 MAC cycles \\
\hline
FENCE & 1 & Pipeline flush \\
\hline
CONFIG\_MVOUT & 1 & Store stride configuration \\
\hline
MVOUT (per 16x16 tile) & 64 & Same as MVIN \\
\hline
\end{tabular}
\caption{RoCC Instruction Latencies}
\end{table}

\subsection{PE Utilization Formula}

\begin{equation}
\text{Utilization} = \frac{\text{actual MACs}}{\text{DIM}^2 \times \text{compute\_cycles}} \times 100\%
\end{equation}

Where:
\begin{itemize}[leftmargin=*]
    \item actual MACs = M x N x K (total multiply-accumulate operations)
    \item DIM = 16 (systolic array dimension)
    \item compute\_cycles = number of cycles the systolic array is active
    \item Theoretical peak MACs/cycle = 256
\end{itemize}

\subsection{Case Study: ResNet-18 Layer-by-Layer}

\begin{table}[H]
\centering
\begin{tabular}{|p{0.25\textwidth}|p{0.15\textwidth}|p{0.15\textwidth}|p{0.15\textwidth}|p{0.15\textwidth}|}
\hline
\textbf{Layer} & \textbf{M} & \textbf{N} & \textbf{K} & \textbf{PE Utilization} \\
\hline
Conv1 (7x7) & 64 & 112x112 & 3x7x7=147 & 88\% \\
\hline
Conv2\_1 (3x3) & 64 & 56x56 & 64x3x3=576 & 94\% \\
\hline
Conv2\_2 (3x3) & 64 & 56x56 & 64x3x3=576 & 94\% \\
\hline
Conv3\_1 (3x3) & 128 & 28x28 & 128x3x3=1152 & 96\% \\
\hline
Conv3\_2 (3x3) & 128 & 28x28 & 128x3x3=1152 & 96\% \\
\hline
Conv4\_1 (3x3) & 256 & 14x14 & 256x3x3=2304 & 97\% \\
\hline
Conv4\_2 (3x3) & 256 & 14x14 & 256x3x3=2304 & 97\% \\
\hline
Conv5\_1 (3x3) & 512 & 7x7 & 512x3x3=4608 & 98\% \\
\hline
Conv5\_2 (3x3) & 512 & 7x7 & 512x3x3=4608 & 98\% \\
\hline
\end{tabular}
\caption{ResNet-18 Convolutional Layers: PE Utilization Estimate}
\end{table}

\begin{warningbox}[Cycle-Accurate Simulation]
Cycle-accurate simulation of Gemmini requires Verilator or FPGA emulation. Spike provides only instruction-accurate simulation (no timing model). The latencies above are analytical estimates based on the Gemmini tile-cycle model.
\end{warningbox}

\section{Performance Benchmarks}

\subsection{Analytical MatMul Throughput}

Hardware model: DIM=16, 1 GHz, 16 B/cycle scratchpad bandwidth.

\begin{table}[H]
\centering
\begin{tabular}{|p{0.2\textwidth}|p{0.2\textwidth}|p{0.2\textwidth}|p{0.2\textwidth}|}
\hline
\textbf{Matrix Size} & \textbf{NumPy host (ms)} & \textbf{Gemmini hw est. (ms)} & \textbf{Speedup} \\
\hline
128 x 128 & 3.038 & 0.014 & 2.09x \\
\hline
256 x 256 & 0.368 & 0.090 & 1.40x \\
\hline
512 x 512 & 1.629 & 0.623 & 1.46x \\
\hline
1024 x 1024 & 9.517 & 4.588 & 1.81x \\
\hline
\end{tabular}
\caption{MatMul Throughput Benchmark}
\end{table}

\subsection{Double-Buffer DMA/Compute Overlap}

\begin{table}[H]
\centering
\begin{tabular}{|p{0.15\textwidth}|p{0.15\textwidth}|p{0.2\textwidth}|p{0.2\textwidth}|p{0.15\textwidth}|}
\hline
\textbf{Matrix Size} & \textbf{Tiles} & \textbf{Seq est. (ms)} & \textbf{Overlap est. (ms)} & \textbf{Benefit} \\
\hline
32 x 32 & 8 & 0.000512 & 0.000384 & 25\% \\
\hline
64 x 64 & 64 & 0.00256 & 0.001536 & 40\% \\
\hline
128 x 128 & 512 & 0.01434 & 0.00819 & 43\% \\
\hline
256 x 256 & 4096 & 0.09011 & 0.06554 & 27\% \\
\hline
\end{tabular}
\caption{Double-Buffer Overlap Benchmark}
\end{table}

\subsection{Numerical Correctness on Spike}

\begin{table}[H]
\centering
\begin{tabular}{|p{0.25\textwidth}|p{0.25\textwidth}|p{0.25\textwidth}|p{0.15\textwidth}|}
\hline
\textbf{Model} & \textbf{Max Abs Error} & \textbf{Max Rel Error} & \textbf{Pass} \\
\hline
ResNet-18 & 1.097e-05 & 9.22e-04 & YES \\
\hline
MobileNetV2 & 5.960e-06 & 7.99e-04 & YES \\
\hline
BERT-tiny & 3.338e-06 & 1.32e-03 & YES \\
\hline
\end{tabular}
\caption{Numerical Correctness at 1e-5 Tolerance}
\end{table}

\section{Usage Guide}

\subsection{Compiling a Model}

\begin{lstlisting}[caption=Basic Compilation Commands, language=bash]
ONNX_MLIR=gemmini_toolchain_build/Release/bin/onnx-mlir

# Emit MLIR for inspection
$ONNX_MLIR --maccel=Gemmini --EmitMLIR model.onnx -o /tmp/model_gemmini

# Count Gemmini runtime calls
grep -c 'om_gemmini_' /tmp/model_gemmini.onnx.mlir

# Compile to RV64 ELF (requires RISC-V toolchain)
$ONNX_MLIR --maccel=Gemmini \
    --mtriple=riscv64-unknown-linux-gnu --mcpu=rocket \
    model.onnx -o model_rv64

# Run on Spike
spike --extension=gemmini --isa=rv64imafdc "$HOME/riscv-gemmini/bin/pk" ./model_rv64
\end{lstlisting}

\subsection{Per-Operator Examples}

\begin{table}[H]
\centering
\begin{tabular}{|p{0.25\textwidth}|p{0.35\textwidth}|p{0.35\textwidth}|}
\hline
\textbf{Operator} & \textbf{Example Command} & \textbf{Expected Symbol} \\
\hline
MatMul f32 & bash gemmini/tests/... & om\_gemmini\_matmul\_f32\_hw \\
\hline
MatMulInteger & bash gemmini/tests/03-.../run\_test.sh & om\_gemmini\_matmulinteger\_i8i8acc32 \\
\hline
QLinearConv & bash gemmini/tests/07-.../run\_test.sh & om\_gemmini\_qlinearconv\_i8\_bias \\
\hline
Conv f32 & bash gemmini/tests/... & om\_gemmini\_conv\_f32\_bias \\
\hline
Resize & bash gemmini/tests/12-resize-f32/run\_test.sh & om\_gemmini\_resize\_nearest\_f32 \\
\hline
Pad & bash gemmini/tests/13-pad-f32/run\_test.sh & om\_gemmini\_pad\_constant\_f32 \\
\hline
Concat & bash gemmini/tests/14-concat-f32/run\_test.sh & om\_gemmini\_concat\_f32 \\
\hline
Slice & bash gemmini/tests/17-slice-f32/run\_test.sh & om\_gemmini\_slice\_f32 \\
\hline
Transpose & bash gemmini/tests/19-transpose-f32/run\_test.sh & om\_gemmini\_transpose\_f32\_hw \\
\hline
Split & bash gemmini/tests/20-split-f32/run\_test.sh & om\_gemmini\_split\_2\_f32 \\
\hline
Sigmoid & bash gemmini/tests/21-sigmoid-f32/run\_test.sh & om\_gemmini\_sigmoid\_f32 \\
\hline
Mul & bash gemmini/tests/22-mul-f32/run\_test.sh & om\_gemmini\_mul\_f32 \\
\hline
\end{tabular}
\caption{Per-Operator Example Commands}
\end{table}

\subsection{Interpreting Emitted MLIR}

\begin{table}[H]
\centering
\begin{tabular}{|p{0.3\textwidth}|p{0.65\textwidth}|}
\hline
\textbf{What you see} & \textbf{Meaning} \\
\hline
funcName = "om\_gemmini\_*" in krnl.call & Runtime-delegate path selected \\
\hline
gemmini\_low.* ops & Direct-RoCC path selected (INT8 MatMul tiles) \\
\hline
llvm.inline\_asm with CUSTOM\_3 & Fully lowered RoCC instruction \\
\hline
Only onnx.* ops remain & Op fell through to standard scalar path \\
\hline
\end{tabular}
\caption{Interpreting Emitted MLIR}
\end{table}

\section{Troubleshooting}

\subsection{Troubleshooting Decision Tree}

\textbf{Tree 1: ``No Gemmini calls in output MLIR''}

\begin{verbatim}
No Gemmini calls
       |
       v
Check build flags -> -DONNX_MLIR_ENABLE_GEMMINI=ON?
       |
       +-- NO -> Reconfigure with -DONNX_MLIR_ENABLE_GEMMINI=ON
       |
       +-- YES
            |
            v
Check ONNX ops -> Are they in the Supported Operators table?
            |
            +-- NO -> Op not yet implemented. Use fallback.
            |
            +-- YES
                 |
                 v
Check dynamic shapes -> Static shapes required for direct path
                 |
                 v
Run with --print-ir-after=onnx-to-gemmini
\end{verbatim}

\textbf{Tree 2: ``Build fails with CMake errors''}

\begin{verbatim}
CMake configuration fails
       |
       v
Run scripts/check_prerequisites.sh --mode <mode>
       |
       +-- Missing gemmini.h -> Set -DGEMMINI_ROOT=/path/to/gemmini
       |
       +-- Missing RISC-V GCC -> Install gcc-riscv64-linux-gnu
       |
       +-- Missing Spike/pk/libgemmini.so
            |
            v
Run scripts/setup_gemmini_simulator.sh --install
\end{verbatim}

\textbf{Tree 3: ``Spike illegal instruction''}

\begin{verbatim}
Spike reports "illegal instruction"
       |
       v
Is --extension=gemmini present?
       |
       +-- NO -> Add --extension=gemmini to Spike command
       |
       +-- YES
            |
            v
Check CUSTOM_X opcode in disassembly
            |
            +-- CUSTOM_0 (0x0b) -> Wrong! Update to CUSTOM_3
            |
            +-- CUSTOM_3 (0x7b) -> Check XCUSTOM_ACC in gemmini_params.h
\end{verbatim}

\subsection{Debugging and Profiling Tools}

\subsubsection{MLIR Debugging Flags}

\begin{lstlisting}[language=bash]
# Print IR after every pass
onnx-mlir-opt --maccel=Gemmini --print-ir-after-all model.mlir

# Print IR before a specific pass
onnx-mlir-opt --maccel=Gemmini --print-ir-before=gemmini-tiling model.mlir

# Print IR after a specific pass
onnx-mlir-opt --maccel=Gemmini --print-ir-after=gemmini-low-to-llvm model.mlir
\end{lstlisting}

\subsubsection{Runtime Verbosity}

\begin{lstlisting}[language=bash]
export OMP_DEBUG=1
export GEMMINI_DEBUG=1
export GEMMINI_VERBOSE=1
\end{lstlisting}

\subsection{Common Issues Table}

% FIX: "Spike illegal instruction on CUSTOM_0" solution now mentions both the
% --extension flag AND the wrong-opcode root cause, so it matches Tree 3.
\begin{longtable}{|p{0.3\textwidth}|p{0.65\textwidth}|}
\hline
\textbf{Issue} & \textbf{Solution} \\
\hline
--maccel=Gemmini not recognized & Rebuild with -DONNX\_MLIR\_ENABLE\_GEMMINI=ON. \\
\hline
Spike illegal instruction on CUSTOM\_0 & Add \texttt{--extension=gemmini} to Spike command; if the opcode in the binary is CUSTOM\_0 (0x0b), the compiler emitted the wrong channel --- update to CUSTOM\_3 (0x7b). See Tree 3. \\
\hline
No Gemmini calls in emitted MLIR & Use memory\_access\_checker.py to see which ops are legal. \\
\hline
Numerical mismatch & Loosen tolerance to 1e-5 or compare with emulator. \\
\hline
Scratchpad range violation & Check tile placements within 16,384 rows (256 KB). \\
\hline
CMake cannot find Gemmini & Set -DGEMMINI\_ROOT=/path/to/gemmini. \\
\hline
Simulator dependencies missing & Run scripts/check\_prerequisites.sh --mode simulator and setup script. \\
\hline
BERT-tiny numerical mismatch & Use decomposed LayerNorm path or tolerance 1e-3. \\
\hline
\end{longtable}

\section{Future Work Roadmap}

\begin{enumerate}[leftmargin=*]
    \item \textbf{Hardware transposer for Conv:} Modify conv lowering to emit config\_ex with OS mode and transpose\_a=true, enabling hardware Im2Col.

    \item \textbf{Im2Col pattern detection:} Add compiler pass that detects Im2Col patterns and replaces them with hardware transposer path.

    \item \textbf{Dynamic shape support via CISC loops:} Expose gemmini\_loop\_ws and gemmini\_loop\_conv\_ws primitives.

    \item \textbf{Native FP16 support:} Extend GemminiTargetInfo.hpp to support inputType = Float(5,10).

    \item \textbf{DMA overlap profiling and tuning:} Profile workloads and tune queue depths and bus widths.

    \item \textbf{Physical hardware validation:} Deploy to Chipyard SoC FPGA and re-run numerical correctness suite.

    \item \textbf{Cycle-accurate profiling:} Use Verilator simulation or bare-metal Spike to access rdcycle.
\end{enumerate}

\section{References and Further Reading}

\subsection{Gemmini Repository}

\begin{itemize}[leftmargin=*]
    \item URL: \url{https://github.com/ucb-bar/gemmini}
    \item Branch: master
    \item Hardware parameters: gemmini/lib/gemmini-rocc-tests/include/gemmini\_params.h
\end{itemize}

\subsection{ONNX-MLIR Repository}

\begin{itemize}[leftmargin=*]
    \item URL: \url{https://github.com/onnx/onnx-mlir}
    \item Gemmini backend location: src/Accelerators/Gemmini/
\end{itemize}

\subsection{Key Papers and Specifications}

\begin{table}[H]
\centering
\begin{tabular}{|p{0.4\textwidth}|p{0.55\textwidth}|}
\hline
\textbf{Reference} & \textbf{Citation} \\
\hline
Gemmini (ASPLOS 2021) & Genc, H., et al. ``Gemmini: Enabling Systematic Deep-Learning Architecture Evaluation via Full-Stack Integration.'' \\
\hline
Systolic Arrays (Kung 1982) & Kung, H. T. ``Why systolic architectures?'' Computer, vol. 15, no. 1, pp. 37-46. \\
\hline
MLIR (CGO 2020) & Lattner, C., et al. ``MLIR: A Compiler Infrastructure for the End of Moore's Law.'' \\
\hline
RISC-V RoCC Interface & RISC-V International, ``RoCC Interface Specification,'' 2019. \\
\hline
ONNX Specification & LF AI \& Data Foundation, ``ONNX Operator Specification v1.17.0,'' 2026. \\
\hline
\end{tabular}
\caption{Key References}
\end{table}

\section{Appendices}

\subsection{Appendix A: Environment Variables}

\begin{longtable}{|p{0.35\textwidth}|p{0.6\textwidth}|}
\hline
\textbf{Variable} & \textbf{Purpose} \\
\hline
GEMMINI\_ROOT & Path to Gemmini installation \\
\hline
RISCV\_INSTALL & Path to RISC-V toolchain and simulator prefix \\
\hline
ONNX\_MLIR\_BUILD\_DIR & Build directory for compiler binaries \\
\hline
LLVM\_INSTALL & Path to LLVM installation \\
\hline
PK & Override proxy kernel path for Spike \\
\hline
SPIKE & Override Spike executable path \\
\hline
IN\_N, IN\_C, IN\_H, IN\_W & Input dimensions for model runners \\
\hline
OMP\_DEBUG & Enable OMTensor debug prints \\
\hline
GEMMINI\_DEBUG & Gemmini runtime debug output \\
\hline
GEMMINI\_VERBOSE & Verbose logging \\
\hline
\end{longtable}

\subsection{Appendix B: CMake Quick Reference}

\begin{lstlisting}[caption=Compile-Only Build, language=bash]
cmake -S . -B gemmini_toolchain_build \
    -DONNX_MLIR_ENABLE_GEMMINI=ON \
    -DONNX_MLIR_GEMMINI_MODE=compile-only \
    -DONNX_MLIR_BUILD_TESTS=ON
\end{lstlisting}

\begin{lstlisting}[caption=Simulator Build, language=bash]
cmake -S . -B gemmini_toolchain_build \
    -DONNX_MLIR_ENABLE_GEMMINI=ON \
    -DONNX_MLIR_GEMMINI_MODE=simulator \
    -DRISCV_INSTALL=/path/to/riscv-gemmini
\end{lstlisting}

\begin{lstlisting}[caption=Hardware Build, language=bash]
cmake -S . -B gemmini_toolchain_build \
    -DONNX_MLIR_ENABLE_GEMMINI=ON \
    -DONNX_MLIR_GEMMINI_MODE=hardware \
    -DRISCV_GCC=riscv64-unknown-elf-gcc
\end{lstlisting}

\subsection{Appendix C: Runtime Function Summary}

\begin{longtable}{|p{0.4\textwidth}|p{0.55\textwidth}|}
\hline
\textbf{Function} & \textbf{Purpose} \\
\hline
om\_gemmini\_tiled\_matmul\_i8i8acc32\_ws & Core int8 MatMul (requires multiples of 16) \\
\hline
om\_gemmini\_quantized\_matmul\_f32\_ws & Quantize f32 to int8 matmul to dequantize \\
\hline
om\_gemmini\_matmul\_f32\_hw & FP32 MatMul \\
\hline
om\_gemmini\_matmul\_f32\_nd\_hw & Batched FP32 MatMul (rank-N x rank-2) \\
\hline
om\_gemmini\_matmul\_f16\_hw & FP16 MatMul via f32 quantized path \\
\hline
om\_gemmini\_conv\_f32, om\_gemmini\_gemm\_f32 & Conv/Gemm using quantized matmul path \\
\hline
om\_gemmini\_convtranspose\_f32 & ConvTranspose (col2im + quantized matmul) \\
\hline
om\_gemmini\_resize\_nearest\_f32, om\_gemmini\_resize\_linear\_f32 & Resize (nearest/linear) \\
\hline
om\_gemmini\_pad\_constant/reflect/edge\_f32 & Pad (constant/reflect/edge) \\
\hline
om\_gemmini\_concat\_f32 & Concat (2-input NCHW) \\
\hline
om\_gemmini\_slice\_f32 & Slice rank-4 with steps and negative indices \\
\hline
om\_gemmini\_transpose\_f32\_hw & Transpose rank-4 (hardware-preferred path) \\
\hline
om\_gemmini\_split\_2/3/4\_f32 & Split 2-4 outputs \\
\hline
om\_gemmini\_sigmoid\_f32 & Sigmoid activation \\
\hline
om\_gemmini\_mul\_f32 & Elementwise multiplication \\
\hline
\end{longtable}

\subsection{Appendix D: Glossary (Extended)}

\begin{longtable}{|p{0.25\textwidth}|p{0.7\textwidth}|}
\hline
\textbf{Term} & \textbf{Definition} \\
\hline
DMA & Direct Memory Access --- data movement between host memory and scratchpad \\
\hline
Scratchpad & On-chip SRAM managed by compiler; 16,384 rows x 16 bytes (256 KB) \\
\hline
Systolic Array & Pipelined compute array for matrix multiplication \\
\hline
Tile & 16x16 block matching systolic array dimensions \\
\hline
Ping-Pong Buffer & Double-buffering technique for DMA/compute overlap \\
\hline
Weight Stationary (WS) & Dataflow where weights are kept in scratchpad across computations \\
\hline
Output Stationary (OS) & Dataflow where partial outputs are kept in accumulator \\
\hline
RoCC & RISC-V Custom Coprocessor interface (CUSTOM\_3 opcode) \\
\hline
Spike & RISC-V architectural simulator \\
\hline
Proxy Kernel (pk) & Lightweight runtime for RISC-V baremetal simulation \\
\hline
Im2Col & Image to Column transformation for convolution \\
\hline
Col2im & Column to Image (scatter-add) for ConvTranspose \\
\hline
CISC Loop & Complex instruction loop for arbitrary matrix sizes \\
\hline
NCHW & Tensor layout: Batch (N), Channel (C), Height (H), Width (W) \\
\hline
LLVM IR & Low-Level Virtual Machine Intermediate Representation \\
\hline
MLIR & Multi-Level Intermediate Representation \\
\hline
PE & Processing Element --- individual MAC unit in systolic array \\
\hline
Arithmetic intensity & Ratio of compute operations to memory accesses (FLOP/byte) \\
\hline
Bank conflict & Multiple memory requests targeting the same bank in the same cycle \\
\hline
Compute-bound & Performance limited by arithmetic throughput \\
\hline
Memory-bound & Performance limited by memory bandwidth \\
\hline
Prefetching & Loading data before it is needed \\
\hline
Gather DMA & DMA reading non-contiguous memory locations \\
\hline
Scatter DMA & DMA writing to non-contiguous memory locations \\
\hline
Strided DMA & DMA with non-unit stride between elements \\
\hline
Accumulator bypass & Direct data path skipping accumulator write \\
\hline
Hardware loop & CISC-style loop instruction without software overhead \\
\hline
Fence instruction & Synchronization primitive for memory ordering \\
\hline
Cache coherence & Protocol keeping caches consistent \\
\hline
Memory barrier & Instruction forcing memory operation ordering \\
\hline
In-flight DMA & DMA transfers issued but not completed \\
\hline
mvin stride & Bytes between rows in DMA load \\
\hline
mvout stride & Bytes between rows in DMA store \\
\hline
Config space & Configuration registers for RoCC commands \\
\hline
RoCC command queue & FIFO buffer for pending RoCC commands \\
\hline
Roofline model & Performance visualization of peak FLOP/s vs bandwidth \\
\hline
\end{longtable}

\subsection{Appendix E: Validation Methodology}

\subsubsection{MLIR Lit Test Structure}

\begin{lstlisting}[caption=Example Lit Test]
// RUN: onnx-mlir-opt --maccel=Gemmini --convert-onnx-to-krnl --canonicalize %s | FileCheck %s

// CHECK-LABEL: func.func @test_conv
// CHECK: %alloc = memref.alloc() : memref<1x64x112x112xf32>
// CHECK: "krnl.call"(%alloc, %arg0, %arg1, %arg2) <{funcName = "om_gemmini_conv_f32_bias"}>
// CHECK-NOT: onnx.Conv
\end{lstlisting}

\subsubsection{Numerical Tolerance Policy}

\begin{table}[H]
\centering
\begin{tabular}{|p{0.25\textwidth}|p{0.2\textwidth}|p{0.2\textwidth}|p{0.3\textwidth}|}
\hline
\textbf{Operation Type} & \textbf{Absolute Tolerance} & \textbf{Relative Tolerance} & \textbf{Reason} \\
\hline
MatMul (f32 quantized) & 1e-5 & 1e-3 & Quantization noise \\
\hline
Conv (f32 quantized) & 1e-5 & 1e-3 & Quantization + im2col error \\
\hline
Elementwise (f32 scalar) & 1e-6 & 1e-4 & Pure scalar, bit-exact \\
\hline
Softmax (f32) & 1e-5 & 1e-3 & Exponential approximation \\
\hline
Resize (nearest) & 0 & 0 & Exact integer mapping \\
\hline
Resize (linear) & 1e-5 & 1e-4 & Floating interpolation \\
\hline
\end{tabular}
\caption{Numerical Tolerance by Operation Type}
\end{table}

\subsubsection{CI/CD Pipeline Diagram}

\begin{verbatim}
CI/CD Pipeline (GitHub Actions)

git push
   |
   v
Ubuntu 22.04 runner (32GB RAM, 8-core)
   |
   +---> Build (compile-only mode)
   |
   +---> MLIR lit tests (41 tests)
   |
   +---> Integration tests (18 tests)
   |
   +---> Example tests (22 tests)
   |
   +---> Model zoo (9 models)
            |
            v
     All tests PASS?
       |        |
      YES       NO
       |        |
       v        v
  Tag release  Fail build
  (v3.0.0)    (email notify)
\end{verbatim}

\begin{notebox}[Validation Environment]
Validation runs on Ubuntu 22.04 LTS with 32GB RAM, 8-core CPU, and 100GB disk space. Spike simulator is installed under ~/riscv-gemmini.
\end{notebox}

\section*{Revision History}

\begin{tabular}{|p{0.2\textwidth}|p{0.2\textwidth}|p{0.2\textwidth}|p{0.35\textwidth}|}
\hline
\textbf{Version} & \textbf{Date} & \textbf{Author} & \textbf{Changes} \\
\hline
1.0 & 2026-05-21 & SST Team & Initial reference manual \\
\hline
2.0 & 2026-05-22 & SST Team & Added FP16, batched MatMul, simulator validation \\
\hline
2.1 & 2026-05-22 & SST Team & Alignment status, test results, runtime verification \\
\hline
3.0 & 2026-06-08 & SST Team & Comprehensive edition with Quick Start, Performance Characterization, troubleshooting, extended glossary \\
\hline
3.1 & \today & SST Team & CORRECTED: Hardware parameters (scratchpad rows: 16,384, 256 KB; accumulator rows: 1,024, 64 KB); Double-buffering slot addresses as relative ranges; RoCC sequence includes FENCE (10 steps, 454 cycles); Lit test table corrected (direct\_matmul\_i8\_rocc.mlir, convtranspose\_gemmini.mlir added; INT8 zero-point tests noted as integration coverage); Common Issues table updated (16,384 rows); Warning-box title corrected (Seven Correctness Invariants); CMake lit target name corrected (check-onnx\_mlir-accelerators-gemmini); CUSTOM\_0 troubleshooting entry updated with opcode fix guidance. \\
\hline
\end{tabular}

\vspace{1cm}

\noindent\textbf{Prepared by:} SST Team, SiliconeSignal Technologies \\
\noindent\textbf{Location:} Meknes, Morocco \\
\noindent\textbf{Document Status:} Final \\
\noindent\textbf{Approval:} Pending

\end{document}
